\section{Conversion des données de la baie CS9 (robot)}
\label{secConversionRobot}

La baie CS9 produit 148 octets et en consomme 124 (cf. document
\texttt{02\_TP/TP\_02}, préparation). Ces octets sont le résultat de
l'encodage, propre au serveur \texttt{uniVALplc}, de l'ensemble des
informations d'état et de commande du robot : ils ne sont exploitables
qu'à travers les blocs \texttt{VAL\_ReadAxesGroup} et
\texttt{VAL\_WriteAxesGroup}.

\subsection{En lecture : \texttt{VAL\_ReadAxesGroup}}

\begin{enumerate}
    \item L'alias \texttt{fromRobotScara}, de type \texttt{T\_FromRobot},
    chevauche la zone image des données produites par la CS9. Comme pour
    le M221, il ne s'agit que d'octets bruts.
    \item \texttt{VAL\_ReadAxesGroup} prend cet objet brut en paramètre
    \texttt{IN} et met à jour, en paramètre \texttt{AxesGroup}, un objet
    \texttt{T\_StaeubliRobot} dont les champs ont un sens explicite.
\end{enumerate}

\texttt{T\_StaeubliRobot} se compose de trois sous-objets :

\begin{center}
\begin{tabular}{|l|p{9.5cm}|}
\hline
\textbf{Champ} & \textbf{Rôle} \\
\hline
\texttt{Status} & État courant du robot, mis à jour par
\texttt{VAL\_ReadAxesGroup} : \texttt{IsMoving}, \texttt{Online},
\texttt{ErrorPending}, \texttt{ActualSpeed}, \texttt{CartesianPos},
\texttt{JointPos}, \texttt{RobotStateMachine}, \texttt{Heartbeat}... \\
\hline
\texttt{Command} & Commandes générales à destination du serveur, écrites
par le programme applicatif puis transmises par
\texttt{VAL\_WriteAxesGroup} : \texttt{OperationModeCmd},
\texttt{OverrideCmd}, \texttt{LifebitPeriod}... \\
\hline
\texttt{CommandInterface} & Interface utilisée en interne par les blocs
\texttt{MC\_} (mouvement, mise sous puissance...) pour transmettre leurs
propres requêtes au serveur. \\
\hline
\end{tabular}
\end{center}

\begin{UPSTIinfor}{Le bit de vie (\texttt{Heartbeat})}
\texttt{Status.Heartbeat} est incrémenté par le serveur \texttt{uniVALplc}
à chaque trame produite. Un programme applicatif rigoureux surveille son
évolution (il doit changer périodiquement, selon
\texttt{Command.LifebitPeriod}) pour détecter une perte de communication
que le seul bit de santé du coupleur ne révélerait pas forcément aussi
rapidement.
\end{UPSTIinfor}

\subsection{Donner des ordres au robot}

Une fois \texttt{T\_StaeubliRobot} disponible et à jour, le programme
applicatif peut agir sur le robot de deux façons :

\begin{itemize}
    \item en renseignant directement certains champs de
    \texttt{Command} (ex. \texttt{OperationModeCmd},
    \texttt{LifebitPeriod}) ;
    \item en appelant des blocs \texttt{MC\_} (\texttt{MC\_GroupPower},
    \texttt{MC\_GroupReset}, \texttt{MC\_MoveLinearAbsolute}...), auxquels
    on transmet l'objet \texttt{AxesGroup} : ces blocs lisent son état
    courant (\texttt{Status}) pour déterminer si une commande peut être
    acceptée, et déposent leur requête dans \texttt{CommandInterface} en
    conséquence.
\end{itemize}

\begin{UPSTIwarning}{Les blocs \texttt{MC\_} ne \og parlent\fg{} jamais directement au réseau}
Un bloc \texttt{MC\_} ne fait que préparer une requête dans l'objet
\texttt{T\_StaeubliRobot} partagé. C'est uniquement au moment où
\texttt{VAL\_WriteAxesGroup} est appelé que cette requête est
effectivement sérialisée et transmise au robot -- voir ci-dessous.
\end{UPSTIwarning}

\subsection{En écriture : \texttt{VAL\_WriteAxesGroup}}

Une fois les ordres déposés dans \texttt{T\_StaeubliRobot}
(\texttt{Command}, \texttt{CommandInterface}), \texttt{VAL\_WriteAxesGroup}
sérialise ces informations vers l'objet brut \texttt{OUT} (de type
\texttt{T\_ToRobot}), qui est ensuite physiquement transmis à la baie CS9
par le coupleur, via l'alias \texttt{toRobotScara}.

\begin{UPSTIidee}{Toujours la même structure en trois temps}
On retrouve, dans l'ordre, les trois étapes suivantes :
\begin{enumerate}
    \item \textbf{Lire} l'objet brut puis le \textbf{convertir} en objet
    sémantique (\texttt{VAL\_ReadAxesGroup}).
    \item \textbf{Donner des ordres} via l'objet sémantique partagé
    (champs de \texttt{Command}, ou blocs \texttt{MC\_}).
    \item \textbf{Convertir} l'objet sémantique en objet brut
    (\texttt{VAL\_WriteAxesGroup}) avant que le coupleur ne
    l'\textbf{écrive} physiquement -- c'est le \textbf{dernier} bloc
    exécuté du cycle, après tous les blocs \texttt{MC\_}.
\end{enumerate}
\end{UPSTIidee}

\UPSTIremarque[Lien avec la préparation du TP]{
C'est précisément ce que demande l'activité « Déclaration des objets
\texttt{T\_FromRobot} et \texttt{T\_ToRobot} » du TP suivant : déclarer,
par chevauchement/alias, les objets \texttt{fromRobotScara} et
\texttt{toRobotScara} qui serviront d'entrée/sortie brutes à
\texttt{VAL\_ReadAxesGroup} et \texttt{VAL\_WriteAxesGroup}.}
